\chapter[Universal Security Bounds]{Universal Security Bounds}\label{ch:sec:universal-security-bounds}

This chapter unifies the smooth and piecewise-smooth theories and establishes a complete classification of security functionals by their regularity. The universal bounds derived here are the culminating results of Part~IV and provide the tools for the security analysis of Chapter~28.


\section{Introduction}\label{ch:sec:introduction-piecewise-smooth-2}
Chapters~\ref{ch:sec:universal-geometric-bounds}--\ref{ch:sec:piecewise-smooth} developed two complementary theories on the WIS
manifold:
\begin{itemize}
\item a \textbf{smooth theory} for differentiable information functionals;
\item a \textbf{piecewise-smooth theory} for ranking-based security functionals.
\end{itemize}
This chapter unifies both perspectives into a single collection of
geometric security bounds expressed in terms of the Universal Optimality
Defect (UOD).
Throughout this chapter,
$\mathcal D(P,Q)=\|F(P)-F(Q)\|^2$
denotes the two-point Universal Optimality Defect.
\subsection{Smooth Security Bounds}\label{ch:thm:18-1}
The results of Chapter~\ref{ch:sec:information-functionals-on-the-wis-manifold} immediately imply the following general
theorem.
\begin{theorem}[Universal Smooth Bound]

Let $S:\mathcal P\rightarrow\mathbb R$ be a \(C^{1}\) functional satisfying
$\|\nabla_gS\|_g\le L.$
Then
\[ \boxed{
|S(P)-S(Q)|
\le
L\sqrt{\mathcal D(P,Q)}.
} \]
If \(S\in C^{2}\) and
$\nabla_gS(P^{\star})=0,$
then locally
\[ \boxed{
|S(P)-S(P^\star)|
\le
C\,\mathcal D(P,P^\star).
} \]
These estimates depend only on the WIS geometry and not on the specific
functional.
\end{theorem}


The results of Chapter~\ref{ch:sec:information-functionals-on-the-wis-manifold} immediately imply the following general
theorem.



\section{A Unified Classification}
Every information or security functional considered in this appendix
falls into one of two categories.
\[
\begin{array}{lll}
\textbf{Functional} & \textbf{Regularity} & \textbf{Applicable Theory} \\ \hline
\text{Shannon entropy} & \text{Smooth} & \text{Chapter~\ref{ch:sec:universal-geometric-bounds}} \\
\text{R\'enyi entropy} & \text{Smooth} & \text{Chapter~\ref{ch:sec:universal-geometric-bounds}} \\
\text{Tsallis entropy} & \text{Smooth} & \text{Chapter~\ref{ch:sec:universal-geometric-bounds}} \\
\text{Collision entropy} & \text{Smooth} & \text{Chapter~\ref{ch:sec:universal-geometric-bounds}} \\
\text{Jensen--Shannon divergence} & \text{Smooth} & \text{Chapters~\ref{ch:sec:universal-geometric-bounds}--\ref{ch:sec:the-average-hessian-operator-and-universal-bregman-theo}} \\
\text{KL divergence} & \text{Smooth/Bregman} & \text{Chapters~\ref{ch:sec:universal-geometric-bounds}--\ref{ch:sec:the-average-hessian-operator-and-universal-bregman-theo}} \\
\text{Min-Entropy} & \text{Piecewise smooth} & \text{Chapter~\ref{ch:sec:piecewise-smooth}} \\
\text{Bayes Vulnerability} & \text{Piecewise smooth} & \text{Chapter~\ref{ch:sec:piecewise-smooth}} \\
\text{Guessing Entropy} & \text{Piecewise smooth} & \text{Chapter~\ref{ch:sec:piecewise-smooth}} \\
\text{Top-\(k\) success probability} & \text{Piecewise smooth} & \text{Chapter~\ref{ch:sec:piecewise-smooth}}
\end{array}
\]
Thus no separate geometric theory is required for each individual
functional.
\subsection{Unified Stability Principle}\label{ch:sec:unified-stability-principle}
The previous results can be summarized by a single principle.
\textbf{Principle.}
Every security functional considered in this appendix admits a geometric
stability estimate controlled by the Universal Optimality Defect.
Specifically,
\begin{itemize}
\item globally for smooth functionals;
\item locally on regularity intervals for piecewise-smooth functionals.
\end{itemize}
The only additional ingredients are regularity assumptions specific to
the functional.


The previous results can be summarized by a single principle.
\textbf{Principle.}
Every security functional considered in this appendix admits a geometric
stability estimate controlled by the Universal Optimality Defect.
Specifically,
\begin{itemize}
\item globally for smooth functionals;
\item locally on regularity intervals for piecewise-smooth functionals.
\end{itemize}
The only additional ingredients are regularity assumptions specific to
the functional.
\subsection{Interpretation}\label{ch:sec:interpretation-piecewise_smooth}
The Universal Optimality Defect plays the role of a canonical geometric
reference quantity.
Rather than deriving independent inequalities for each entropy or
security measure, the WIS framework separates the analysis into two
independent components:
\begin{enumerate}
\item the geometry of the posterior manifold, encoded by the UOD;
\item the analytical regularity of the functional.
\end{enumerate}
This separation isolates the geometric contribution from the
functional-specific one.


\subsection{Scope and Limitations}\label{ch:sec:scope-and-limitations}
The results established in this appendix should be interpreted as
\textbf{stability estimates} rather than optimal inequalities.
In particular,
\begin{itemize}
\item the constants generally depend on the region under consideration;
\item piecewise bounds do not automatically extend across ranking transitions;
\item global constants may fail to exist near the boundary of the probability simplex.
\end{itemize}
\paragraph{Running example: browser fingerprinting.}
The min-entropy $H_\infty(P)=-\log\max_i P_i$ of the browser fingerprint distribution (Section~\ref{ch:sec:browser-fingerprinting-and-anonymity}) is the canonical piecewise-smooth functional: on the region $\{P\in\Delta_n^\circ:P_1>P_2>\dots>P_n\}$ it is smooth with $H_\infty(P)=-\log P_1$, but at a boundary $P_i=P_j$ the identity of the maximiser changes abruptly and the gradient $\nabla H_\infty$ jumps.  Theorem~\ref{ch:thm:18-1} guarantees that on each interval the Lipschitz bound $|H_\infty(P)-H_\infty(Q)|\le L\sqrt{\mathcal D(P,Q)}$ holds with the same $L$, so a small perturbation of the fingerprint distribution cannot produce an arbitrarily large change in identifiability.

These limitations arise from the analytical properties of the
functionals rather than from the WIS geometry itself.


The results established in this appendix should be interpreted as
\textbf{stability estimates} rather than optimal inequalities.
In particular,
\begin{itemize}
\item the constants generally depend on the region under consideration;
\item piecewise bounds do not automatically extend across ranking transitions;
\item global constants may fail to exist near the boundary of the probability simplex.
\end{itemize}
\paragraph{Running example: browser fingerprinting.}
The min-entropy $H_\infty(P)=-\log\max_i P_i$ of the browser fingerprint distribution (Section~\ref{ch:sec:browser-fingerprinting-and-anonymity}) is the canonical piecewise-smooth functional: on the region $\{P\in\Delta_n^\circ:P_1>P_2>\dots>P_n\}$ it is smooth with $H_\infty(P)=-\log P_1$, but at a boundary $P_i=P_j$ the identity of the maximiser changes abruptly and the gradient $\nabla H_\infty$ jumps.  Theorem~\ref{ch:thm:18-1} guarantees that on each interval the Lipschitz bound $|H_\infty(P)-H_\infty(Q)|\le L\sqrt{\mathcal D(P,Q)}$ holds with the same $L$, so a small perturbation of the fingerprint distribution cannot produce an arbitrarily large change in identifiability.

These limitations arise from the analytical properties of the
functionals rather than from the WIS geometry itself.
